% Preamble template for Tufte-style lecture notes.
% Do not compile this file directly; compile the main notes file instead.
\documentclass{tufte-handout}

\usepackage{ntheorem}
\usepackage{graphicx}
\usepackage{amsmath}
\usepackage{amssymb}
\usepackage{hyperref}
\usepackage{epigraph}
\usepackage{booktabs}
\theoremstyle{break}
% \usepackage[
% bibencoding=utf8,% .bib file encoding
% maxbibnames=3, % otherwise et al
% minbibnames=1, % otherwise et al
% backend=biber,%
% sortlocale=en_US,%
% style=apa,% or authoryear
% % apabackref=false, % backreferences
% natbib=true,% for citet/citep, but this is for backward compatibility
% uniquename=false,%
% url=true,%
% sortcites=false,
% doi=true,%
% eprint=true%
% ]{biblatex}
% \addbibresource{lecture_note_bib.bib}

\hypersetup{
  colorlinks,
  urlcolor = blue,
  pdfauthor={Paul Goldsmith-Pinkham}
  pdfkeywords={econometrics}
  pdftitle={Lecture Notes for Applied Empirical Methods}
  pdfpagemode=UseNone
}
\newtheorem{ruleN}{Rule}
\newtheorem{thmN}{Theorem}
\newtheorem{assN}{Assumption}
\newtheorem{defN}{Definition}
\newtheorem{exmp}{Example}
\newtheorem{cmt}{Comment}
\newtheorem{proof}{Proof}

\newcommand{\continuation}{??}
\newtheorem*{excont}{Example \continuation}
\newenvironment{continueexample}[1]
 {\renewcommand{\continuation}{\ref{#1}}\excont[continued]}
 {\endexcont}
\newcommand{\bY}{\mathbf{Y}}
\newcommand{\bX}{\mathbf{X}}
\newcommand{\bD}{\mathbf{D}}
\newcommand{\E}{\mathbb{E}}

\newcommand\independent{\protect\mathpalette{\protect\independenT}{\perp}}
\def\independenT#1#2{\mathrel{\rlap{$#1#2$}\mkern2mu{#1#2}}}
\DeclareMathOperator{\Supp}{Supp}

\usepackage{cleveref}
\crefname{appsec}{appendix}{appendices}
\crefname{appsubsec}{appendix}{appendices}
\crefname{assumption}{assumption}{assumptions}
\crefname{equation}{equation}{equations}
\crefname{exmp}{example}{examples}
\crefname{assN}{assumption}{assumptions}
\crefname{cmt}{comment}{comments}
\crefname{defN}{definition}{definitions}

\usepackage[nolist]{acronym}
\begin{acronym}
  \acro{CI}{confidence interval}%
  \acro{OLS}{ordinary least squares}%
  \acro{CLT}{central limit theorem}%
  \acro{IV}{instrumental variables}%
  \acro{ATE}{average treatment effect}%
  \acro{RCT}{randomized control trial}%
  \acro{SUTVA}{stable unit treatment value assignment}
  \acro{VAM}{value-added model}%
  \acro{LAN}{locally asymptotically normal}%
  \acro{DiD}{difference-in-differences}%
  \acro{OVB}{omitted variables bias}
  \acro{FWL}{Frisch-Waugh-Lovell}
  \acro{DAG}{directed acyclic graph}
  \acro{PO}{potential outcomes}
\end{acronym}

\def\inprobHIGH{\,{\buildrel p \over \rightarrow}\,} 
\def\inprob{\,{\inprobHIGH}\,} 
\def\indistHIGH{\,{\buildrel d \over \rightarrow}\,} 
\def\indist{\,{\indistHIGH}\,}

\usepackage[many]{tcolorbox} 
\definecolor{main}{HTML}{3F6EA7}      % primary accent
\definecolor{sub}{HTML}{F2F7FD}       % cool background
\definecolor{sub2}{HTML}{FFF5E8}      % warm background
\definecolor{mainline}{HTML}{D3DFEE}  % soft frame for information boxes
\definecolor{warn}{HTML}{C7771B}      % accent for caution boxes
\definecolor{warnline}{HTML}{EAD9C0}  % soft frame for caution boxes

\tcbset{
    enhanced,
    breakable,
    colback = white,
    boxrule = 0.5pt,
    arc = 2pt,
    left = 9pt,
    right = 9pt,
    top = 7pt,
    bottom = 7pt,
    before skip = 0.3cm,
    after skip = 0.45cm
}                           % setting global options for tcolorbox

\newtcolorbox{boxD}{
    colback = sub,
    colframe = mainline,
    borderline west = {2.2pt}{0pt}{main},
    borderline north = {0.6pt}{0pt}{main!25},
    borderline south = {0.6pt}{0pt}{main!15}
}


\newtcolorbox{boxF}{
    colback = sub2,
    colframe = warnline,
    borderline west = {2.2pt}{0pt}{warn},
    borderline north = {0.6pt}{0pt}{warn!25},
    borderline south = {0.6pt}{0pt}{warn!15}
}

\usepackage{colortbl}


\usepackage{tikz}
\usepackage{verbatim}
\usetikzlibrary{positioning}
\usetikzlibrary{snakes}
\usetikzlibrary{calc}
\usetikzlibrary{arrows}
\usetikzlibrary{decorations.markings}
\usetikzlibrary{shapes.misc}
\usetikzlibrary{matrix,shapes,arrows,fit,tikzmark}




\title{S181M116 Derivatives and Alternative Investments: Mathematical Foundations}
\author{Swapnil Singh}
\date{}

\hypersetup{
  pdftitle={S181M116 Derivatives and Alternative Investments: Mathematical Foundations},
  pdfauthor={Swapnil Singh},
  pdfkeywords={derivatives, mathematical foundations, algebra, probability, calculus, discounting, stochastic processes}
}

\setlength{\headheight}{26pt}

\newcommand{\R}{\mathbb{R}}
\newcommand{\Prob}{\mathbb{P}}
\newcommand{\Var}{\operatorname{Var}}
\newcommand{\Cov}{\operatorname{Cov}}
\newcommand{\Corr}{\operatorname{Corr}}
\newcommand{\NPV}{\mathrm{NPV}}
\newcommand{\PV}{\mathrm{PV}}
\newcommand{\FV}{\mathrm{FV}}
\newcommand{\St}{S_T}
\newcommand{\Szero}{S_0}
\newcommand{\Fzero}{F_0}
\newcommand{\K}{K}
\newcommand{\dd}{\,d}

\begin{document}
\maketitle

\begin{abstract}
This note reviews the mathematical foundations needed for the course. It is not
a general mathematics textbook. It is a practical toolkit for derivatives and
alternative investments: algebra, discounting, payoff functions, no-arbitrage
logic, probability, normal and lognormal distributions, calculus, Taylor
approximations, optimization, portfolio variance, regression beta, binomial
trees, and the first ideas behind Brownian motion and It\^o's lemma. The note is
example-heavy because the course is mathematically intensive and the formulas are
useful only when students can interpret them economically.
\end{abstract}

\begin{boxD}
\textbf{How to use this note.} Read Sections 1--6 before the first class if the
notation is unfamiliar. Return to Sections 7--12 before the lectures on hedging,
binomial trees, Black--Scholes--Merton, Greeks, volatility, and real options.
The goal is working fluency: you should be able to compute, draw, explain, and
check each object.
\end{boxD}

\section{Why Mathematics Matters in This Course}

Derivatives are contracts whose cash flows depend on future variables. This
means that almost every important question has a mathematical form:
\[
    \text{contract value today}
    =
    \text{value of uncertain future cash flows}.
\]
The course repeatedly uses four ideas.
\begin{enumerate}
    \item \textbf{Payoffs are functions.} A forward payoff, option payoff, swap
    payment, and carried-interest payoff are all functions of future variables.
    \item \textbf{Time matters.} A euro today and a euro next year are not the
    same object. Discounting converts future cash flows into present values.
    \item \textbf{Risk matters.} Future cash flows are random. We need
    probability, expectation, variance, covariance, and distributional thinking.
    \item \textbf{Replication matters.} If two strategies generate the same
    future cash flows, they should have the same current value. This is the
    no-arbitrage principle.
\end{enumerate}

\begin{boxF}
\textbf{Central habit.} Do not memorize a formula without asking three questions:
What is being measured? What are the units? Under what assumptions is the formula
valid?
\end{boxF}

\subsection{A Map of the Toolkit}

\begin{tabular}{@{}p{0.25\linewidth}p{0.65\linewidth}@{}}
\toprule
Course topic & Mathematical tools used \\
\midrule
Forwards and futures & Algebra, compounding, discounting, no-arbitrage
replication, linear payoff functions \\
Hedging with futures & Covariance, correlation, regression beta, variance
minimization \\
Interest rates and swaps & Present value, discount factors, zero rates, forward
rates, sums of cash flows \\
Options & Positive-part functions, convexity, inequalities, put-call parity,
piecewise functions \\
Binomial trees & Weighted averages, replication, backward induction, risk-neutral
probabilities \\
Black--Scholes--Merton & Calculus, partial derivatives, Taylor expansion,
Brownian motion, It\^o's lemma \\
Greeks and risk & First and second derivatives, scenario approximation,
portfolio aggregation \\
Alternative investments & Net present value, IRR, option-like flexibility,
discounted cash-flow logic \\
\bottomrule
\end{tabular}

\section{Notation and Algebra}

\subsection{Variables, Parameters, and Subscripts}

A \emph{variable} can change. A \emph{parameter} is treated as fixed within the
calculation. In derivatives notation:
\[
    S_0 = \text{spot price today}, \qquad S_T = \text{spot price at maturity},
\]
\[
    K = \text{strike or delivery price}, \qquad T = \text{time to maturity}.
\]
The subscript $0$ usually means today. The subscript $T$ means at maturity.

\begin{exmp}[Reading notation]
If a European call has payoff
\[
    (S_T-K)^+,
\]
then $S_T$ is random before maturity, while $K$ is fixed by the option contract.
The plus sign does not mean ``add something.'' It means positive part:
\[
    x^+ = \max(x,0).
\]
\end{exmp}

\subsection{Units}

Financial formulas are easier when units are explicit. If $S_T$ and $K$ are both
in dollars per share, then $S_T-K$ is also dollars per share. If the contract is
for $N$ shares, the total payoff is
\[
    N(S_T-K).
\]

\begin{exmp}[Units in a forward payoff]
A forward contract is written on 5,000 barrels of oil. The delivery price is
$K=74$ dollars per barrel. At maturity the spot price is $S_T=79$ dollars per
barrel. The long forward payoff per barrel is
\[
    S_T-K = 79-74 = 5.
\]
The total payoff is
\[
    5{,}000 \times 5 = \$25{,}000.
\]
The unit check is
\[
    \text{barrels} \times \frac{\text{dollars}}{\text{barrel}}
    =
    \text{dollars}.
\]
\end{exmp}

\subsection{Fractions, Percentages, and Basis Points}

A percentage rate must be converted into decimal form before it is used in a
formula:
\[
    5\% = 0.05, \qquad 12.5\% = 0.125.
\]
One basis point is one-hundredth of one percentage point:
\[
    1\text{ bp}=0.01\%=0.0001.
\]
Therefore 75 basis points is $0.75\%$ or $0.0075$.

\begin{exmp}[Basis points]
An interest rate increases from $3.20\%$ to $3.65\%$. The change is
\[
    3.65\%-3.20\%=0.45\%=45\text{ bp}.
\]
In decimal form the rate increases from $0.0320$ to $0.0365$.
\end{exmp}

\subsection{Powers and Exponents}

The expression $x^n$ means $x$ multiplied by itself $n$ times. The expression
$e^x$ uses the exponential function, where $e\approx 2.71828$. Exponentials are
central because continuously compounded interest grows as $e^{rT}$.

Useful rules:
\[
    e^a e^b = e^{a+b}, \qquad \frac{e^a}{e^b}=e^{a-b}, \qquad (e^a)^b=e^{ab}.
\]
For any positive number $A$,
\[
    \ln(A)
\]
is the natural logarithm, the inverse of the exponential:
\[
    \ln(e^x)=x, \qquad e^{\ln(A)}=A.
\]

\begin{exmp}[Solving with logs]
Suppose a price grows from 100 to 110 over one year under continuous compounding.
Find the continuously compounded return $r$:
\[
    110=100e^r.
\]
Divide by 100:
\[
    1.10=e^r.
\]
Take logs:
\[
    r=\ln(1.10)\approx 0.0953.
\]
The continuously compounded return is $9.53\%$.
\end{exmp}

\subsection{Summation Notation}

The symbol $\sum$ means add terms:
\[
    \sum_{t=1}^{T} CF_t
    =
    CF_1+CF_2+\cdots+CF_T.
\]
For discounted cash flows,
\[
    \PV=\sum_{t=1}^{T}\frac{CF_t}{(1+r)^t}.
\]

\begin{exmp}[Present value as a sum]
A project pays 100 at the end of year 1, 120 at the end of year 2, and 150 at
the end of year 3. If the annual discount rate is 5\%, then
\[
    \PV = \frac{100}{1.05}+\frac{120}{1.05^2}+\frac{150}{1.05^3}
    \approx 333.18.
\]
\end{exmp}

\section{Time Value of Money}

\subsection{Future Value and Present Value}

If money earns interest, one unit today grows into more than one unit tomorrow.
With annual compounding at rate $r$,
\[
    \FV = \PV(1+r)^T.
\]
Solving for present value gives
\[
    \PV = \frac{\FV}{(1+r)^T}.
\]

\begin{defN}[Discount factor]
The discount factor $D(0,T)$ is the price today of one unit of currency paid at
time $T$ with no default risk. Under annual compounding,
\[
    D(0,T)=\frac{1}{(1+r)^T}.
\]
Under continuous compounding,
\[
    D(0,T)=e^{-rT}.
\]
\end{defN}

\begin{exmp}[Discounting one cash flow]
You will receive 1,000 in two years. The annual interest rate is 4\%. The present
value is
\[
    \PV=\frac{1{,}000}{1.04^2}\approx 924.56.
\]
Under continuous compounding at 4\%, the present value is
\[
    \PV=1{,}000e^{-0.04\times 2}\approx 923.12.
\]
\end{exmp}

\subsection{Compounding Conventions}

With $m$ compounding periods per year, the future value of one unit invested for
$T$ years at nominal annual rate $r$ is
\[
    \left(1+\frac{r}{m}\right)^{mT}.
\]
As $m$ becomes very large, this approaches continuous compounding:
\[
    \lim_{m\to\infty}\left(1+\frac{r}{m}\right)^{mT}=e^{rT}.
\]

\begin{exmp}[Annual, quarterly, and continuous compounding]
Invest 100 for two years at a quoted annual rate of 6\%.
\[
\begin{aligned}
    \text{Annual compounding:} \quad
    &100(1.06)^2 = 112.36,\\
    \text{Quarterly compounding:} \quad
    &100\left(1+\frac{0.06}{4}\right)^8 \approx 112.65,\\
    \text{Continuous compounding:} \quad
    &100e^{0.06\times 2}\approx 112.75.
\end{aligned}
\]
More frequent compounding produces slightly more interest for the same nominal
rate.
\end{exmp}

\subsection{Converting Rates}

If $R_c$ is a continuously compounded rate and $R_m$ is a rate compounded $m$
times per year, the same one-year accumulation requires
\[
    e^{R_c}=\left(1+\frac{R_m}{m}\right)^m.
\]
Therefore
\[
    R_c = m\ln\left(1+\frac{R_m}{m}\right),
    \qquad
    R_m = m\left(e^{R_c/m}-1\right).
\]

\begin{exmp}[Convert an annual rate to continuous compounding]
A bank quotes 8\% with annual compounding. The equivalent continuously
compounded rate is
\[
    R_c=\ln(1.08)\approx 0.07696.
\]
So 8\% annual compounding is equivalent to about 7.696\% continuous compounding.
\end{exmp}

\subsection{Zero Rates, Discount Factors, and Forward Rates}

A zero-coupon bond pays one cash flow at maturity. If the continuously
compounded zero rate for maturity $T$ is $R(0,T)$, then the discount factor is
\[
    P(0,T)=e^{-R(0,T)T}.
\]
If we know discount factors, we can recover zero rates:
\[
    R(0,T)=-\frac{1}{T}\ln P(0,T).
\]

The forward rate from $T_1$ to $T_2$ is the rate agreed today for borrowing or
lending between those future dates. With continuous compounding,
\[
    f(T_1,T_2)
    =
    \frac{R(0,T_2)T_2-R(0,T_1)T_1}{T_2-T_1}.
\]

\begin{exmp}[Forward rate from zero rates]
The one-year continuously compounded zero rate is 3\% and the two-year zero rate
is 4\%. The one-year forward rate starting one year from now is
\[
    f(1,2)=\frac{0.04\times 2-0.03\times 1}{2-1}=0.05.
\]
The implied forward rate is 5\%.
\end{exmp}

\subsection{Net Present Value}

Net present value subtracts the initial cost from the present value of future
cash flows:
\[
    \NPV = -I_0 + \sum_{t=1}^{T}\frac{CF_t}{(1+r)^t}.
\]
The financial interpretation is simple: if $\NPV>0$, the project is worth more
than it costs using discount rate $r$.

\begin{exmp}[Project NPV]
A real asset investment costs 500 today and is expected to generate 180 at the
end of each of the next three years. At a discount rate of 7\%,
\[
    \NPV=-500+\frac{180}{1.07}+\frac{180}{1.07^2}+\frac{180}{1.07^3}
    \approx -27.63.
\]
At this discount rate, the project destroys value unless it has additional
strategic value, such as an option to expand later.
\end{exmp}

\section{Functions and Payoffs}

\subsection{Functions}

A function maps an input into an output. In finance, the input is often a future
price and the output is a payoff:
\[
    \text{payoff}=g(S_T).
\]
The same function can be represented by a formula, a table, or a graph.

\begin{exmp}[A linear payoff function]
A long forward with delivery price $K=100$ has payoff
\[
    g(S_T)=S_T-100.
\]
\begin{tabular}{@{}rrrrr@{}}
\toprule
$S_T$ & 70 & 90 & 100 & 130 \\
\midrule
$g(S_T)$ & -30 & -10 & 0 & 30 \\
\bottomrule
\end{tabular}
The payoff increases one-for-one with the terminal spot price.
\end{exmp}

\subsection{Positive-Part Notation}

The positive part of $x$ is
\[
    x^+ = \max(x,0)
    =
    \begin{cases}
    x, & x>0,\\
    0, & x\leq 0.
    \end{cases}
\]
This notation is used constantly for options:
\[
    \text{Call payoff}=(S_T-K)^+,
    \qquad
    \text{Put payoff}=(K-S_T)^+.
\]

\begin{exmp}[Call option payoff]
A call has strike $K=50$. Its payoff is
\[
    (S_T-50)^+.
\]
\begin{tabular}{@{}rrrrrr@{}}
\toprule
$S_T$ & 35 & 45 & 50 & 60 & 80 \\
\midrule
Call payoff & 0 & 0 & 0 & 10 & 30 \\
\bottomrule
\end{tabular}
The option is valuable at maturity only when the underlying price exceeds the
strike.
\end{exmp}

\begin{exmp}[Put option payoff]
A put has strike $K=50$. Its payoff is
\[
    (50-S_T)^+.
\]
\begin{tabular}{@{}rrrrrr@{}}
\toprule
$S_T$ & 35 & 45 & 50 & 60 & 80 \\
\midrule
Put payoff & 15 & 5 & 0 & 0 & 0 \\
\bottomrule
\end{tabular}
The put is valuable at maturity when the underlying price is below the strike.
\end{exmp}

\subsection{Payoff Versus Profit}

Payoff is the cash flow at maturity. Profit subtracts the initial price of the
position, usually adjusted for financing. Ignoring financing, a call bought for
premium $c_0$ has profit
\[
    (S_T-K)^+ - c_0.
\]

\begin{exmp}[Option payoff and profit]
A trader buys a call with strike 100 for premium 7. Ignoring interest on the
premium:
\[
    \text{profit}=(S_T-100)^+-7.
\]
\begin{tabular}{@{}rrrrrr@{}}
\toprule
$S_T$ & 80 & 95 & 100 & 107 & 120 \\
\midrule
Payoff & 0 & 0 & 0 & 7 & 20 \\
Profit & -7 & -7 & -7 & 0 & 13 \\
\bottomrule
\end{tabular}
The break-even terminal price is 107.
\end{exmp}

\subsection{Piecewise Functions}

Many derivative payoffs are piecewise: the formula changes depending on the
state. A call payoff can be written as
\[
    (S_T-K)^+
    =
    \begin{cases}
    0, & S_T\leq K,\\
    S_T-K, & S_T>K.
    \end{cases}
\]

\begin{exmp}[Covered call payoff]
A covered call consists of buying one share and selling one call with strike
$K$. Ignoring the option premium, the terminal payoff is
\[
    S_T-(S_T-K)^+.
\]
If $S_T\leq K$, the call expires worthless and the payoff is $S_T$. If $S_T>K$,
the short call loses $S_T-K$, so the total payoff is
\[
    S_T-(S_T-K)=K.
\]
Thus
\[
    S_T-(S_T-K)^+
    =
    \begin{cases}
    S_T, & S_T\leq K,\\
    K, & S_T>K.
    \end{cases}
\]
The covered call keeps downside stock exposure but caps the upside.
\end{exmp}

\subsection{Convexity}

A function is convex when the line segment between two points on the graph lies
above the graph. Intuitively, a convex payoff benefits from movement in either
direction. Options have convex payoffs.

For a differentiable function, convexity corresponds to a nonnegative second
derivative:
\[
    g''(x)\geq 0.
\]
Linear functions are both convex and concave because their second derivative is
zero.

\begin{exmp}[Why a call is convex]
The call payoff is flat when $S_T<K$ and then rises with slope one when $S_T>K$.
The slope jumps upward from 0 to 1 at $K$. That upward change in slope is the
source of convexity.
\end{exmp}

\begin{boxF}
\textbf{Financial intuition.} Forwards give linear exposure. Options give
nonlinear exposure. Nonlinearity is why option pricing requires more machinery
than forward pricing.
\end{boxF}

\section{No-Arbitrage Logic}

\subsection{Law of One Price}

The law of one price says that two assets with exactly the same future cash
flows in every state should have the same price today. If not, a trader can buy
the cheap asset and sell the expensive asset.

\begin{defN}[Arbitrage]
An arbitrage is a trading strategy with no initial cost or a negative initial
cost, no possibility of future loss, and a positive probability of future gain.
\end{defN}

Most textbook pricing arguments use a stronger practical version: if two
self-financing strategies have the same future payoff, their current costs must
be equal.

\subsection{Forward Price for a Non-Dividend-Paying Asset}

Consider a stock with current price $S_0$ and no dividends before time $T$.
Compare two strategies:
\begin{enumerate}
    \item Buy the stock today and finance the purchase by borrowing $S_0$.
    \item Enter a forward contract to buy the stock at time $T$ for delivery
    price $F_0$.
\end{enumerate}
At time $T$, both strategies produce ownership of one share. The first strategy
requires repayment $S_0e^{rT}$. The second strategy requires payment $F_0$.
No-arbitrage implies
\[
    F_0 = S_0e^{rT}.
\]

\begin{exmp}[Cash-and-carry arbitrage]
Suppose $S_0=100$, $r=5\%$, $T=1$, and the market forward price is 108. The
no-arbitrage forward price is
\[
    100e^{0.05}\approx 105.13.
\]
The forward is too expensive. A trader can:
\begin{enumerate}
    \item borrow 100,
    \item buy the stock,
    \item short the forward at 108.
\end{enumerate}
At maturity, deliver the stock into the forward, receive 108, and repay
$100e^{0.05}=105.13$. The locked-in profit is
\[
    108-105.13=2.87.
\]
\end{exmp}

\begin{exmp}[Reverse cash-and-carry arbitrage]
Suppose $S_0=100$, $r=5\%$, $T=1$, and the market forward price is 102. The
forward is too cheap relative to $105.13$. A trader can:
\begin{enumerate}
    \item short sell the stock and receive 100,
    \item invest 100 at the risk-free rate,
    \item enter a long forward at 102.
\end{enumerate}
At maturity, the investment grows to 105.13. Use 102 to buy the stock through the
forward and return it to close the short sale. The locked-in profit is
\[
    105.13-102=3.13.
\]
\end{exmp}

\subsection{Option Bounds}

No-arbitrage also gives inequalities. A European call on a non-dividend-paying
stock cannot be worth more than the stock:
\[
    c \leq S_0.
\]
It cannot be worth less than
\[
    S_0-Ke^{-rT},
\]
unless this number is negative. Therefore
\[
    c \geq \max(S_0-Ke^{-rT},0).
\]

\begin{exmp}[Detecting a call-price violation]
Let $S_0=50$, $K=45$, $r=4\%$, and $T=0.5$. The lower bound is
\[
    50-45e^{-0.04\times 0.5}
    \approx 50-44.11
    =
    5.89.
\]
If the call trades for 4.80, it is too cheap relative to the bound. The quoted
price violates no-arbitrage under the assumptions.
\end{exmp}

\subsection{Put-Call Parity}

For European options on a non-dividend-paying stock, put-call parity is
\[
    c + Ke^{-rT} = p + S_0.
\]
The left side is a call plus a risk-free bond that pays $K$ at maturity. The
right side is a put plus the stock. At maturity both sides pay
\[
    \max(S_T,K).
\]

\begin{exmp}[Put-call parity]
Suppose $S_0=100$, $K=100$, $r=5\%$, $T=1$, and a European call price is
$c=10.45$. Put-call parity implies
\[
    p = c + Ke^{-rT} - S_0
    = 10.45 + 100e^{-0.05} - 100
    \approx 5.57.
\]
If the put trades at 7.00, the put is expensive relative to the call, stock, and
bond combination.
\end{exmp}

\begin{boxD}
\textbf{Method.} No-arbitrage problems usually become easy when you build a
payoff table. Write the payoff of each strategy in each future state. If two
columns are identical, their current costs should match.
\end{boxD}

\section{Probability}

\subsection{Random Variables}

A random variable is a numerical outcome of an uncertain event. Before maturity,
$S_T$ is a random variable. After maturity, it is observed.

\begin{exmp}[Random payoff]
A stock will be either 120 or 80 in one year. A call with strike 100 has payoff
\[
    (S_T-100)^+.
\]
\begin{tabular}{@{}lrr@{}}
\toprule
State & $S_T$ & Call payoff \\
\midrule
Up & 120 & 20 \\
Down & 80 & 0 \\
\bottomrule
\end{tabular}
The payoff is random because it depends on the future state.
\end{exmp}

\subsection{Expected Value}

The expected value is the probability-weighted average outcome. If $X$ takes
values $x_1,\ldots,x_n$ with probabilities $p_1,\ldots,p_n$, then
\[
    \E[X]=\sum_{i=1}^{n}p_i x_i.
\]

\begin{exmp}[Expected stock price]
A stock will be 120 with probability 0.6 and 80 with probability 0.4. Then
\[
    \E[S_T]=0.6(120)+0.4(80)=104.
\]
The expected price is not necessarily a price at which the stock can trade. It
is a probability-weighted average.
\end{exmp}

\begin{exmp}[Expected option payoff]
With the same states, a call with strike 100 pays 20 in the up state and 0 in
the down state. Its expected payoff is
\[
    \E[(S_T-100)^+]=0.6(20)+0.4(0)=12.
\]
This is not automatically the option price. Pricing also requires discounting
and risk adjustment or replication.
\end{exmp}

\subsection{Variance and Standard Deviation}

Variance measures average squared deviation from the mean:
\[
    \Var(X)=\E[(X-\E[X])^2].
\]
Standard deviation is the square root of variance:
\[
    \sigma_X=\sqrt{\Var(X)}.
\]
In finance, standard deviation is often called volatility.

\begin{exmp}[Variance in a two-state model]
Let $S_T$ be 120 with probability 0.6 and 80 with probability 0.4. We already
computed $\E[S_T]=104$. The variance is
\[
\begin{aligned}
    \Var(S_T)
    &=
    0.6(120-104)^2+0.4(80-104)^2\\
    &=
    0.6(16)^2+0.4(-24)^2\\
    &=
    384.
\end{aligned}
\]
The standard deviation is
\[
    \sqrt{384}\approx 19.60.
\]
\end{exmp}

\subsection{Covariance and Correlation}

Covariance measures whether two variables move together:
\[
    \Cov(X,Y)=\E[(X-\E[X])(Y-\E[Y])].
\]
Correlation rescales covariance:
\[
    \Corr(X,Y)=\frac{\Cov(X,Y)}{\sigma_X\sigma_Y}.
\]
Correlation always lies between $-1$ and $1$.

\begin{exmp}[Interpreting correlation]
If oil-price changes and airline stock returns have negative correlation, oil
price increases tend to coincide with airline stock losses. This does not mean
oil prices are the only cause of airline returns. Correlation is a co-movement
measure, not a complete causal statement.
\end{exmp}

\begin{exmp}[Covariance calculation]
Suppose two returns have the following equally likely outcomes.
\[
\begin{array}{@{}lrr@{}}
\toprule
State & R_A & R_B \\
\midrule
1 & 0.10 & 0.06 \\
2 & 0.02 & 0.01 \\
3 & -0.04 & -0.02 \\
\bottomrule
\end{array}
\]
The means are
\[
\begin{aligned}
    \E[R_A]&=\frac{0.10+0.02-0.04}{3}=0.0267,\\
    \E[R_B]&=\frac{0.06+0.01-0.02}{3}=0.0167.
\end{aligned}
\]
The covariance is
\[
\begin{aligned}
    \Cov(R_A,R_B)
    &=
    \frac{1}{3}
    \left[
    (0.10-0.0267)(0.06-0.0167)\right.\\
    &\quad
    +(0.02-0.0267)(0.01-0.0167)\\
    &\quad
    \left.
    +(-0.04-0.0267)(-0.02-0.0167)
    \right]\\
    &\approx 0.00189.
\end{aligned}
\]
The covariance is positive because the returns move in the same direction across
states.
\end{exmp}

\subsection{Conditional Probability and Conditional Expectation}

Conditional probability updates probabilities after information is observed.
The notation
\[
    \Prob(A\mid B)
\]
means the probability of event $A$ given that event $B$ has occurred.

Conditional expectation is the expected value after conditioning on information:
\[
    \E[X\mid \mathcal{I}].
\]
In finance, $\mathcal{I}$ often represents current market information.

\begin{exmp}[Conditional expectation]
A commodity price can be high or low next month. If a storage report is normal,
the expected price is 70. If the report shows a shortage, the expected price is
82. Before seeing the report, the expectation averages over both possibilities.
After seeing the report, the conditional expectation uses the relevant
information set.
\end{exmp}

\subsection{Independent Events and Independent Random Variables}

Events $A$ and $B$ are independent if
\[
    \Prob(A\cap B)=\Prob(A)\Prob(B).
\]
Random variables are independent if knowing one does not change the distribution
of the other.

\begin{boxF}
\textbf{Warning.} Zero correlation is weaker than independence. If two variables
are independent, their covariance is zero. But covariance can be zero even when
the variables are not independent.
\end{boxF}

\section{Normal and Lognormal Distributions}

\subsection{Normal Distribution}

A normal random variable has the familiar bell-shaped distribution. We write
\[
    X\sim N(\mu,\sigma^2)
\]
to mean that $X$ is normally distributed with mean $\mu$ and variance
$\sigma^2$.

The standard normal distribution is
\[
    Z\sim N(0,1).
\]
Any normal variable can be standardized:
\[
    Z=\frac{X-\mu}{\sigma}.
\]

\begin{exmp}[Standardization]
Suppose daily portfolio returns are approximately normal with mean 0.05\% and
standard deviation 1.20\%. What is the standardized value of a return of
$-2.35\%$?
\[
    z=\frac{-0.0235-0.0005}{0.0120}=-2.
\]
The return is two standard deviations below the mean.
\end{exmp}

\subsection{Why Normal Distributions Appear in Finance}

Normal distributions are useful because sums of many small shocks are often
approximately normal. This is the central limit theorem intuition. In the
Black--Scholes--Merton model, continuously compounded returns over a fixed
horizon are normally distributed.

\begin{boxD}
\textbf{Interpretation.} A normal distribution can take negative values. That is
acceptable for returns, but not for stock prices. This is one reason the model
uses normal log returns and lognormal prices.
\end{boxD}

\subsection{Lognormal Distribution}

A positive random variable $S$ is lognormal if its logarithm is normal:
\[
    \ln S \sim N(\mu,\sigma^2).
\]
Equivalently,
\[
    S=e^X
\]
where $X$ is normal.

Lognormal variables are always positive. This makes them natural for stock
prices in the Black--Scholes--Merton model.

\begin{exmp}[From normal log return to lognormal price]
Suppose
\[
    \ln(S_T/S_0)\sim N(0.04,0.20^2).
\]
Then $S_T/S_0$ is lognormal, and $S_T$ is always positive. A one-standard
deviation high log return is
\[
    0.04+0.20=0.24,
\]
which corresponds to a gross return
\[
    e^{0.24}\approx 1.271.
\]
That is a 27.1\% simple return.
\end{exmp}

\subsection{Simple Returns and Log Returns}

The simple return from $0$ to $T$ is
\[
    R=\frac{S_T-S_0}{S_0}.
\]
The log return is
\[
    r=\ln\left(\frac{S_T}{S_0}\right).
\]
They are related by
\[
    r=\ln(1+R), \qquad R=e^r-1.
\]
For small returns, $r\approx R$.

\begin{exmp}[Simple and log return]
A stock rises from 50 to 55. The simple return is
\[
    R=\frac{55-50}{50}=0.10.
\]
The log return is
\[
    r=\ln(55/50)=\ln(1.10)\approx 0.0953.
\]
For a 10\% simple return, the log return is 9.53\%.
\end{exmp}

\begin{exmp}[Why log returns add]
A stock rises from 100 to 110 and then falls from 110 to 99. The two simple
returns are $10\%$ and $-10\%$, but the total simple return is $-1\%$.

The log returns are
\[
    \ln(110/100)=\ln(1.10),
    \qquad
    \ln(99/110)=\ln(0.90).
\]
The total log return is
\[
    \ln(1.10)+\ln(0.90)=\ln(0.99),
\]
which is exactly the log return from 100 to 99.
\end{exmp}

\section{Calculus for Finance}

\subsection{Derivatives as Slopes}

The derivative of $f(x)$ is the instantaneous rate of change of $f$ with respect
to $x$:
\[
    f'(x)=\frac{\dd f}{\dd x}.
\]
In finance, derivatives measure sensitivities. The option Greek Delta is
\[
    \Delta=\frac{\partial f}{\partial S},
\]
the sensitivity of option value to the underlying price.

\begin{exmp}[Derivative of a linear payoff]
For a long forward payoff
\[
    g(S_T)=S_T-K,
\]
the derivative with respect to $S_T$ is
\[
    g'(S_T)=1.
\]
The payoff changes one-for-one with the terminal spot price.
\end{exmp}

\begin{exmp}[Derivative of a bond price with respect to yield]
A zero-coupon bond paying 100 at time $T$ has price
\[
    B(y)=100e^{-yT}.
\]
The derivative with respect to yield is
\[
    \frac{\dd B}{\dd y}=-T100e^{-yT}=-TB(y).
\]
The negative sign means that bond prices fall when yields rise.
\end{exmp}

\subsection{Common Derivative Rules}

\[
    \frac{\dd}{\dd x}x^n = nx^{n-1},
    \qquad
    \frac{\dd}{\dd x}e^x=e^x,
    \qquad
    \frac{\dd}{\dd x}\ln x=\frac{1}{x}.
\]
Product rule:
\[
    \frac{\dd}{\dd x}[f(x)g(x)]
    =
    f'(x)g(x)+f(x)g'(x).
\]
Chain rule:
\[
    \frac{\dd}{\dd x}f(g(x))=f'(g(x))g'(x).
\]

\begin{exmp}[Chain rule in discounting]
Let
\[
    B(y)=e^{-yT}.
\]
Set $u=-yT$. Then $B=e^u$ and $\dd u/\dd y=-T$. The chain rule gives
\[
    \frac{\dd B}{\dd y}=e^u(-T)=-Te^{-yT}.
\]
\end{exmp}

\subsection{Partial Derivatives}

When a function has several inputs, a partial derivative changes one input while
holding the others fixed. If an option value is
\[
    f=f(S,t,\sigma,r),
\]
then
\[
    \frac{\partial f}{\partial S}
\]
changes $S$ while holding $t,\sigma,r$ fixed.

\begin{exmp}[Partial derivatives]
Let
\[
    f(S,t)=S^2e^{-rt}.
\]
Then
\[
    \frac{\partial f}{\partial S}=2Se^{-rt},
    \qquad
    \frac{\partial f}{\partial t}=-rS^2e^{-rt}.
\]
The first derivative measures sensitivity to the price. The second measures
sensitivity to time, holding $S$ fixed.
\end{exmp}

\subsection{Taylor Approximation}

Taylor approximation replaces a nonlinear function with a local polynomial.
First order:
\[
    f(x+\Delta x)\approx f(x)+f'(x)\Delta x.
\]
Second order:
\[
    f(x+\Delta x)
    \approx
    f(x)+f'(x)\Delta x+\frac{1}{2}f''(x)(\Delta x)^2.
\]

For an option value $f(S,t)$, the standard Greek approximation is
\[
    \Delta f
    \approx
    \Delta \Delta S
    +\frac{1}{2}\Gamma(\Delta S)^2
    +\Theta \Delta t.
\]
Here the first $\Delta$ on the right side is the Greek Delta, while $\Delta S$
is the change in the underlying price.

\begin{exmp}[Delta-Gamma approximation]
An option has $\Delta=0.60$ and $\Gamma=0.04$. The stock price rises by 2. The
Delta-Gamma approximation to the option-value change is
\[
    0.60(2)+\frac{1}{2}(0.04)(2)^2
    =
    1.20+0.08
    =
    1.28.
\]
The Gamma term matters because the option value is curved in the stock price.
\end{exmp}

\subsection{Optimization}

Optimization chooses a variable to maximize or minimize an objective. In one
dimension, an interior optimum usually satisfies the first-order condition
\[
    f'(x)=0.
\]
For a minimum, the second derivative should be positive:
\[
    f''(x)>0.
\]

\begin{exmp}[Minimum-variance hedge ratio]
Let $\Delta S$ be the change in the spot price and $\Delta F$ the change in the
futures price. A short hedge using hedge ratio $h$ has change
\[
    \Delta V=\Delta S-h\Delta F.
\]
The variance is
\[
    \Var(\Delta V)
    =
    \sigma_S^2+h^2\sigma_F^2-2h\Cov(\Delta S,\Delta F).
\]
Differentiate with respect to $h$:
\[
    \frac{\dd}{\dd h}\Var(\Delta V)
    =
    2h\sigma_F^2-2\Cov(\Delta S,\Delta F).
\]
Set this equal to zero:
\[
    h^*
    =
    \frac{\Cov(\Delta S,\Delta F)}{\sigma_F^2}
    =
    \rho\frac{\sigma_S}{\sigma_F}.
\]
This is the minimum-variance hedge ratio.
\end{exmp}

\section{Matrices and Portfolio Risk}

\subsection{Vectors}

A vector is an ordered list of numbers. A portfolio with weights
\[
    w=
    \begin{pmatrix}
    w_1\\
    w_2\\
    \vdots\\
    w_n
    \end{pmatrix}
\]
and asset returns
\[
    R=
    \begin{pmatrix}
    R_1\\
    R_2\\
    \vdots\\
    R_n
    \end{pmatrix}
\]
has portfolio return
\[
    R_p=w^\top R=\sum_{i=1}^{n}w_i R_i.
\]

\begin{exmp}[Portfolio return]
A portfolio invests 40\% in asset A and 60\% in asset B. If returns are
$R_A=8\%$ and $R_B=3\%$, then
\[
    R_p=0.4(0.08)+0.6(0.03)=0.05.
\]
The portfolio return is 5\%.
\end{exmp}

\subsection{Covariance Matrix}

The covariance matrix collects variances and covariances:
\[
    \Sigma=
    \begin{pmatrix}
    \sigma_1^2 & \sigma_{12}\\
    \sigma_{12} & \sigma_2^2
    \end{pmatrix}.
\]
For portfolio weights $w$, portfolio variance is
\[
    \sigma_p^2=w^\top\Sigma w.
\]
For two assets,
\[
    \sigma_p^2
    =
    w_1^2\sigma_1^2+w_2^2\sigma_2^2+2w_1w_2\sigma_{12}.
\]

\begin{exmp}[Two-asset portfolio variance]
Asset A has volatility 20\%, asset B has volatility 10\%, and their correlation
is 0.30. The portfolio weights are 50\% and 50\%. The covariance is
\[
    \sigma_{AB}=\rho\sigma_A\sigma_B=0.30(0.20)(0.10)=0.006.
\]
The portfolio variance is
\[
    0.5^2(0.20)^2+0.5^2(0.10)^2+2(0.5)(0.5)(0.006)
    =
    0.0155.
\]
The portfolio volatility is
\[
    \sqrt{0.0155}\approx 0.1245.
\]
So the portfolio volatility is about 12.45\%.
\end{exmp}

\subsection{Diversification}

Diversification works when assets are not perfectly correlated. If two assets
have correlation less than 1, the portfolio volatility can be below the weighted
average of individual volatilities.

\begin{boxD}
\textbf{Derivative connection.} A hedge is a portfolio designed so that one
position offsets risk in another. The effectiveness of the hedge depends on
covariance, not only on the size of the two positions.
\end{boxD}

\section{Statistics Used in the Course}

\subsection{Sample Mean and Sample Variance}

Given observations $x_1,\ldots,x_n$, the sample mean is
\[
    \bar{x}=\frac{1}{n}\sum_{i=1}^{n}x_i.
\]
The sample variance is
\[
    s^2=\frac{1}{n-1}\sum_{i=1}^{n}(x_i-\bar{x})^2.
\]
The denominator $n-1$ adjusts for the fact that the mean is estimated from the
same data.

\begin{exmp}[Sample volatility]
Suppose five daily returns are
\[
    0.4\%,\quad -0.6\%,\quad 1.0\%,\quad 0.2\%,\quad -0.5\%.
\]
In decimals:
\[
    0.004,\quad -0.006,\quad 0.010,\quad 0.002,\quad -0.005.
\]
The sample mean is
\[
    \bar{x}=0.001.
\]
The sample variance is
\[
\begin{aligned}
    s^2
    &=
    \frac{1}{4}\left[
    (0.004-0.001)^2+(-0.006-0.001)^2\right.\\
    &\quad
    +(0.010-0.001)^2+(0.002-0.001)^2\\
    &\quad
    \left.+(-0.005-0.001)^2
    \right]\\
    &=
    0.000044.
\end{aligned}
\]
The sample standard deviation is
\[
    s=\sqrt{0.000044}\approx 0.00663,
\]
or 0.663\% per day.
\end{exmp}

\subsection{Annualizing Volatility}

If daily returns are independent with constant variance, annual variance is
approximately daily variance times the number of trading days. With 252 trading
days,
\[
    \sigma_{\text{annual}} \approx \sigma_{\text{daily}}\sqrt{252}.
\]

\begin{exmp}[Annualized volatility]
If daily volatility is 1.2\%, annualized volatility is
\[
    0.012\sqrt{252}\approx 0.1905.
\]
So the annualized volatility is about 19.05\%.
\end{exmp}

\begin{boxF}
\textbf{Assumption.} The square-root-of-time rule relies on weak dependence and
roughly constant variance. It can be misleading during crises, for illiquid
assets, or when volatility changes sharply.
\end{boxF}

\subsection{Regression Beta}

A simple regression of asset returns on market returns is
\[
    R_i = \alpha + \beta R_m + \varepsilon.
\]
The slope is
\[
    \beta=\frac{\Cov(R_i,R_m)}{\Var(R_m)}.
\]
Beta measures market exposure.

\begin{exmp}[Equity index hedge]
A portfolio is worth USD 10 million and has beta 1.2 relative to an index. The
index futures price is 5,000 and each futures contract is for 50 times the index.
The notional per futures contract is
\[
    5{,}000\times 50=\$250{,}000.
\]
The number of contracts for a short hedge is
\[
    N^*=\beta\frac{V_A}{V_F}
    =
    1.2\frac{10{,}000{,}000}{250{,}000}
    =
    48.
\]
The portfolio manager shorts 48 futures contracts to reduce market exposure.
\end{exmp}

\section{Binomial Trees and Risk-Neutral Valuation}

\subsection{One-Step Binomial Model}

In a one-step binomial model, the stock price moves from $S_0$ to either
\[
    S_u=S_0u
    \qquad \text{or} \qquad
    S_d=S_0d,
\]
where $u>d$.

For a derivative with payoffs $f_u$ and $f_d$, replication uses $\Delta$ shares
and a bond position $B$:
\[
    \Delta S_0u + Be^{rT}=f_u,
\]
\[
    \Delta S_0d + Be^{rT}=f_d.
\]
Subtract the second equation from the first:
\[
    \Delta=\frac{f_u-f_d}{S_0u-S_0d}.
\]

\begin{exmp}[One-step call by replication]
Let $S_0=100$, $u=1.2$, $d=0.8$, $K=100$, $r=5\%$, and $T=1$. The stock prices
are 120 and 80. The call payoffs are
\[
    f_u=20,\qquad f_d=0.
\]
The hedge ratio is
\[
    \Delta=\frac{20-0}{120-80}=0.5.
\]
Solve for the bond position using the down state:
\[
    0.5(80)+Be^{0.05}=0.
\]
Thus
\[
    B=-40e^{-0.05}\approx -38.05.
\]
The option value is
\[
    f_0=\Delta S_0+B=0.5(100)-38.05=11.95.
\]
\end{exmp}

\subsection{Risk-Neutral Probability}

The same price can be written as a discounted expected payoff under the
risk-neutral probability
\[
    p=\frac{e^{rT}-d}{u-d}.
\]
Then
\[
    f_0=e^{-rT}[pf_u+(1-p)f_d].
\]

\begin{exmp}[One-step call by risk-neutral probability]
Using the previous example,
\[
    p=\frac{e^{0.05}-0.8}{1.2-0.8}\approx 0.6282.
\]
The option price is
\[
    f_0=e^{-0.05}[0.6282(20)+(1-0.6282)(0)]
    \approx 11.95.
\]
This matches the replication price.
\end{exmp}

\begin{boxD}
\textbf{Interpretation.} Risk-neutral probabilities are pricing weights, not
necessarily real-world probabilities. They are chosen so that discounted asset
prices are consistent with no-arbitrage.
\end{boxD}

\subsection{Backward Induction}

In a multi-step tree, start at maturity and move backward one node at a time:
\[
    f=e^{-r\Delta t}[pf_u+(1-p)f_d].
\]
For American options, compare continuation value with immediate exercise value
at each node.

\begin{exmp}[American put early exercise logic]
At a node, suppose an American put has strike 50 and the stock price is 42.
Immediate exercise pays
\[
    50-42=8.
\]
Suppose the discounted continuation value from holding the option is 7.40. The
American value at the node is
\[
    \max(8,7.40)=8.
\]
Exercise is optimal at that node in the tree.
\end{exmp}

\section{Stochastic Processes}

\subsection{From Random Variables to Processes}

A stochastic process is a collection of random variables indexed by time:
\[
    \{S_t: t\geq 0\}.
\]
Instead of one uncertain future value, we model an uncertain path over time.

\begin{exmp}[Daily price process]
A stock price observed every day can be written as
\[
    S_0,S_1,S_2,\ldots.
\]
Before the future days arrive, $S_1,S_2,\ldots$ are random variables. After the
days pass, they become observed values.
\end{exmp}

\subsection{Markov Property}

A process is Markov if the future depends on the present state, not on the full
history:
\[
    \Prob(S_{t+\Delta t}\leq x \mid S_t,S_{t-1},S_{t-2},\ldots)
    =
    \Prob(S_{t+\Delta t}\leq x \mid S_t).
\]
This is a simplifying assumption, not a universal truth.

\begin{boxF}
\textbf{Modeling warning.} The Markov property helps make pricing tractable. It
does not mean real markets have no memory. Volatility clustering, liquidity
stress, and institutional constraints often create path dependence.
\end{boxF}

\subsection{Brownian Motion}

A standard Brownian motion, or Wiener process, is usually written $z_t$. It has:
\begin{enumerate}
    \item $z_0=0$,
    \item independent increments,
    \item normally distributed increments,
    \item continuous paths.
\end{enumerate}
Over a small time interval $\Delta t$,
\[
    \Delta z = \varepsilon\sqrt{\Delta t},
    \qquad
    \varepsilon\sim N(0,1).
\]
Thus
\[
    \E[\Delta z]=0,
    \qquad
    \Var(\Delta z)=\Delta t.
\]

\begin{exmp}[Brownian increment]
If $\Delta t=1/252$ year, then
\[
    \sqrt{\Delta t}=\sqrt{1/252}\approx 0.0630.
\]
A one-standard-deviation Brownian increment over one trading day is about
0.0630 in annual time units.
\end{exmp}

\subsection{Geometric Brownian Motion}

The Black--Scholes--Merton model assumes the stock price follows
\[
    dS = \mu S\,dt+\sigma S\,dz.
\]
This is geometric Brownian motion. The term $\mu S\,dt$ is drift. The term
$\sigma S\,dz$ is random shock.

Dividing by $S$ gives
\[
    \frac{dS}{S}=\mu\,dt+\sigma\,dz.
\]
So $\mu$ is expected continuously compounded return per unit time and $\sigma$
is volatility per square-root unit of time.

\begin{exmp}[Interpreting the GBM shock]
Suppose annual volatility is $\sigma=20\%$ and $\Delta t=1/252$. A one-standard
deviation daily price shock is approximately
\[
    \sigma\sqrt{\Delta t}=0.20\sqrt{1/252}\approx 0.0126.
\]
That is about 1.26\% of the stock price.
\end{exmp}

\subsection{Quadratic Variation Intuition}

In ordinary calculus, terms like $(dt)^2$ are negligible. In stochastic
calculus, Brownian increments satisfy
\[
    (dz)^2=dt,
\]
while
\[
    dt\,dz=0,
    \qquad
    (dt)^2=0.
\]
The informal reason is that $dz$ is of order $\sqrt{dt}$, so $(dz)^2$ is of
order $dt$.

\begin{boxD}
\textbf{Key idea.} The rule $(dz)^2=dt$ is why It\^o's lemma has an extra
second-derivative term. That term becomes central in the Black--Scholes--Merton
PDE.
\end{boxD}

\subsection{It\^o's Lemma}

If
\[
    dx=a(x,t)\,dt+b(x,t)\,dz,
\]
and $G=G(x,t)$, then It\^o's lemma says
\[
    dG=
    \left(
    \frac{\partial G}{\partial x}a
    +\frac{\partial G}{\partial t}
    +\frac{1}{2}\frac{\partial^2G}{\partial x^2}b^2
    \right)dt
    +
    \frac{\partial G}{\partial x}b\,dz.
\]

\begin{exmp}[Log of a GBM stock price]
Let
\[
    dS=\mu S\,dt+\sigma S\,dz
\]
and define
\[
    G(S)=\ln S.
\]
The derivatives are
\[
    G_S=\frac{1}{S},
    \qquad
    G_{SS}=-\frac{1}{S^2},
    \qquad
    G_t=0.
\]
It\^o's lemma gives
\[
\begin{aligned}
    d\ln S
    &=
    \left[
    \frac{1}{S}(\mu S)
    +\frac{1}{2}\left(-\frac{1}{S^2}\right)(\sigma^2S^2)
    \right]dt
    +
    \frac{1}{S}(\sigma S)\,dz\\
    &=
    \left(\mu-\frac{1}{2}\sigma^2\right)dt+\sigma\,dz.
\end{aligned}
\]
The log price has drift $\mu-\frac{1}{2}\sigma^2$, not $\mu$.
\end{exmp}

\section{Black--Scholes--Merton Building Blocks}

\subsection{Delta Hedging}

If an option has value $f(S,t)$, its small change is approximately
\[
    df \approx f_S\,dS + f_t\,dt + \frac{1}{2}f_{SS}(dS)^2.
\]
The option's exposure to the stock over a small move is $f_S$, called Delta.
A portfolio that is long one option and short $\Delta=f_S$ shares has stock
exposure approximately canceled over a very short interval.

\begin{exmp}[Delta-neutral position]
An option has Delta 0.65. A trader is long 100 options, each written on one
share. The stock exposure is approximately
\[
    100\times 0.65=65
\]
shares. To make the position delta-neutral, the trader shorts 65 shares.
\end{exmp}

\subsection{Greeks as Sensitivities}

\begin{tabular}{@{}lll@{}}
\toprule
Greek & Formula & Interpretation \\
\midrule
Delta & $\Delta=\partial f/\partial S$ & sensitivity to underlying price \\
Gamma & $\Gamma=\partial^2 f/\partial S^2$ & curvature in underlying price \\
Theta & $\Theta=\partial f/\partial t$ & sensitivity to calendar time \\
Vega & $\nu=\partial f/\partial \sigma$ & sensitivity to volatility \\
Rho & $\rho=\partial f/\partial r$ & sensitivity to interest rate \\
\bottomrule
\end{tabular}

\begin{exmp}[Greek-based scenario]
A portfolio has Delta $-20{,}000$, Gamma 500, and Vega 75,000. The underlying
price increases by 1 and implied volatility increases by 0.01. Ignoring time and
interest-rate effects, the approximate value change is
\[
\begin{aligned}
    -20{,}000(1)+\frac{1}{2}(500)(1)^2+75{,}000(0.01)
    &=
    -20{,}000+250+750
    \\
    &=
    -19{,}000.
\end{aligned}
\]
\end{exmp}

\subsection{Risk-Neutral Valuation}

Risk-neutral valuation says that, under no-arbitrage, derivative prices can be
computed as discounted expected payoffs under a special probability measure:
\[
    f_0=e^{-rT}\E^Q[g(S_T)].
\]
The superscript $Q$ means the expectation is taken under risk-neutral
probabilities.

For a stock under the risk-neutral measure,
\[
    dS=rS\,dt+\sigma S\,dz^Q
\]
in the simplest no-dividend Black--Scholes--Merton model. The expected return
changes from $\mu$ to $r$, but volatility remains $\sigma$.

\begin{boxF}
\textbf{Common confusion.} Risk-neutral valuation does not say investors are
risk-neutral. It says derivative prices can be computed as if expected returns
were the risk-free rate after risk has been correctly embedded in pricing
probabilities.
\end{boxF}

\section{Alternative Investment Mathematics}

\subsection{Internal Rate of Return}

The internal rate of return is the discount rate that sets net present value to
zero:
\[
    0=\sum_{t=0}^{T}\frac{CF_t}{(1+\mathrm{IRR})^t}.
\]
The initial investment is usually a negative cash flow.

\begin{exmp}[IRR for a two-year investment]
An investment costs 100 today, pays 30 after one year, and pays 90 after two
years. The IRR solves
\[
    -100+\frac{30}{1+r}+\frac{90}{(1+r)^2}=0.
\]
Let $x=1+r$. Then
\[
    -100+\frac{30}{x}+\frac{90}{x^2}=0.
\]
Multiply by $x^2$:
\[
    -100x^2+30x+90=0.
\]
Equivalently,
\[
    100x^2-30x-90=0.
\]
The positive solution is
\[
    x=\frac{30+\sqrt{30^2+4(100)(90)}}{200}\approx 1.1105.
\]
Thus
\[
    r\approx 0.1105.
\]
The IRR is about 11.05\%.
\end{exmp}

\begin{boxF}
\textbf{IRR warning.} IRR can be misleading when cash flows change signs more
than once, when projects have different scale, or when interim cash flows cannot
be reinvested at the IRR. NPV is usually the cleaner value measure.
\end{boxF}

\subsection{Money Multiples}

Private equity often reports multiples:
\[
    \mathrm{TVPI}
    =
    \frac{\text{distributions}+\text{residual value}}{\text{paid-in capital}}.
\]
DPI uses distributions only. RVPI uses residual value only.

\begin{exmp}[Private equity multiples]
A fund has paid-in capital of 200, cumulative distributions of 90, and remaining
net asset value of 170. Then
\[
    \mathrm{DPI}=\frac{90}{200}=0.45,
    \qquad
    \mathrm{RVPI}=\frac{170}{200}=0.85,
\]
and
\[
    \mathrm{TVPI}=0.45+0.85=1.30.
\]
The fund has generated total value equal to 1.30 times paid-in capital, but much
of it remains unrealized.
\end{exmp}

\subsection{Real Options}

A real option is managerial flexibility with option-like payoff. Examples:
\begin{itemize}
    \item option to delay investment,
    \item option to expand,
    \item option to abandon,
    \item option to switch inputs or outputs,
    \item option to stage financing.
\end{itemize}

\begin{exmp}[Option to abandon]
A project has uncertain continuation value next year. If continued, it is worth
either 150 or 60. The firm can abandon the project and sell assets for 80. The
value next year is
\[
    \max(V,80).
\]
In the high state this is 150. In the low state this is 80, not 60. The
abandonment option creates a floor, just like a put option creates downside
protection.
\end{exmp}

\begin{exmp}[Option to delay]
A firm can invest 100 today in a project whose value is uncertain next year. If
the firm waits, it can invest only if the project value is high. Waiting has an
opportunity cost because cash flows are delayed, but it also has option value
because the firm avoids investing in bad states. The trade-off is the same logic
as exercising versus holding an option.
\end{exmp}

\section{Common Mistakes and Diagnostic Checks}

\subsection{Mistake 1: Mixing Percent and Decimal Units}

Using 5 instead of 0.05 in an interest-rate formula changes the calculation from
5\% to 500\%.

\begin{exmp}[Unit error]
The correct one-year growth factor at 5\% is
\[
    e^{0.05}=1.0513.
\]
The incorrect calculation
\[
    e^5=148.41
\]
would imply a 14,741\% return.
\end{exmp}

\subsection{Mistake 2: Treating Expected Payoff as Price}

The expected payoff is not generally the price. A derivative price must account
for discounting, risk, and replication.

\begin{exmp}[Expected payoff is not enough]
A risky payoff has expected value 100 in one year. If it is very risky, its
price may be far below $100e^{-rT}$. If it can be perfectly replicated, its price
is determined by the replication cost, not by a subjective expected payoff.
\end{exmp}

\subsection{Mistake 3: Ignoring Contract Size}

Futures and options quotes are often quoted per unit, but contracts cover many
units.

\begin{exmp}[Contract multiplier]
An index option quote changes by 3.20 index points. If the contract multiplier is
100, the dollar change per contract is
\[
    3.20\times 100=\$320.
\]
For 25 contracts, the total change is
\[
    25\times 320=\$8{,}000.
\]
\end{exmp}

\subsection{Mistake 4: Confusing Real-World and Risk-Neutral Probabilities}

Real-world probabilities describe statistical beliefs about future states.
Risk-neutral probabilities are pricing weights that make no-arbitrage valuation
work.

\begin{exmp}[Two different probabilities]
In a one-step tree, the real-world probability of an up move may be 0.55. But
the risk-neutral probability used for pricing may be 0.48. The derivative price
uses 0.48 in the no-arbitrage formula, not because 0.48 is the true forecast,
but because it is the pricing weight consistent with the stock and bond prices.
\end{exmp}

\section{Worked Review Problems}

\begin{exmp}[Forward price with income]
A stock is worth 80 today and will pay a known dividend with present value 2
before maturity. The continuously compounded risk-free rate is 4\% and maturity
is one year. The forward price is
\[
    F_0=(S_0-I)e^{rT}=(80-2)e^{0.04}\approx 81.18.
\]
The dividend lowers the forward price because the forward buyer does not receive
the dividend before delivery.
\end{exmp}

\begin{exmp}[Put-call parity arbitrage direction]
Let $S_0=40$, $K=40$, $r=5\%$, $T=1$, $c=4.50$, and $p=1.00$. Put-call parity
requires
\[
    c+Ke^{-rT}=p+S_0.
\]
The left side is
\[
    4.50+40e^{-0.05}\approx 42.55.
\]
The right side is
\[
    1.00+40=41.00.
\]
The left side is expensive. Sell the left side and buy the right side:
short the call, borrow the present value of 40 by shorting the bond, buy the put,
and buy the stock. This locks in the price difference under the assumptions of
European exercise, no dividends, and frictionless trading.
\end{exmp}

\begin{exmp}[Minimum-variance hedge]
A firm has exposure to a spot commodity. The monthly standard deviation of spot
price changes is 3.0, the monthly standard deviation of futures price changes is
2.5, and the correlation is 0.80. The optimal hedge ratio is
\[
    h^*=\rho\frac{\sigma_S}{\sigma_F}
    =
    0.80\frac{3.0}{2.5}
    =
    0.96.
\]
If the firm wants to hedge 100,000 units and each futures contract is for 10,000
units, the number of contracts is
\[
    N^*=0.96\frac{100{,}000}{10{,}000}=9.6.
\]
In practice the firm must choose 9 or 10 contracts, then consider transaction
costs, liquidity, and basis risk.
\end{exmp}

\begin{exmp}[Discounted expected payoff under risk-neutral probability]
A derivative pays 30 in the up state and 0 in the down state after one year.
The risk-neutral probability of the up state is 0.40 and the risk-free rate is
5\% continuously compounded. The price is
\[
    e^{-0.05}[0.40(30)+0.60(0)]
    =
    11.41.
\]
\end{exmp}

\begin{exmp}[Delta-Gamma-Vega approximation]
A trader owns an option portfolio with Delta 8,000, Gamma $-300$, and Vega
45,000. The underlying falls by 1.5 and implied volatility rises by 0.02.
Ignoring Theta and Rho, the approximate change is
\[
\begin{aligned}
    8{,}000(-1.5)
    +\frac{1}{2}(-300)(-1.5)^2
    +45{,}000(0.02)
    &=
    -12{,}000-337.5+900
    \\
    &=
    -11{,}437.5.
\end{aligned}
\]
\end{exmp}

\section{Checklist for Students}

After reading this note, you should be able to:
\begin{itemize}
    \item convert percentages, decimals, and basis points;
    \item use annual, periodic, and continuous compounding;
    \item discount a single cash flow and a stream of cash flows;
    \item compute zero rates, discount factors, and simple forward rates;
    \item read payoff notation such as $(S_T-K)^+$;
    \item distinguish payoff, profit, and value;
    \item build payoff tables for forwards, calls, puts, and combinations;
    \item state the law of one price and use it in no-arbitrage arguments;
    \item compute expectation, variance, covariance, and correlation;
    \item distinguish simple returns and log returns;
    \item explain why lognormal prices are positive;
    \item compute basic derivatives and partial derivatives;
    \item use Taylor approximations to interpret Greeks;
    \item derive the minimum-variance hedge ratio;
    \item compute portfolio variance from variances and covariance;
    \item annualize volatility using the square-root-of-time rule;
    \item price a one-step binomial option by replication and risk-neutral
    probability;
    \item explain the intuition behind Brownian motion, GBM, and It\^o's lemma;
    \item compute IRR and money multiples for simple alternative investment
    examples.
\end{itemize}

\section{Practice Questions}

\begin{enumerate}
    \item Convert 125 basis points into a percentage and a decimal.
    \item A cash flow of 10,000 arrives in 18 months. The continuously compounded
    risk-free rate is 3.5\%. Compute its present value.
    \item An investment grows from 75 to 90 over two years. Compute the simple
    two-year return and the continuously compounded annual return.
    \item A long forward has delivery price 62. Compute the payoff for
    $S_T=45,62,70,88$.
    \item A call has strike 100 and premium 6. Compute payoff and profit for
    $S_T=80,100,106,125$, ignoring financing.
    \item A put has strike 50 and premium 3. Compute payoff and profit for
    $S_T=35,47,50,60$, ignoring financing.
    \item Explain in words why a call option payoff is convex.
    \item A non-dividend-paying stock has $S_0=120$, $r=4\%$, and $T=0.75$.
    Compute the no-arbitrage forward price using continuous compounding.
    \item If the market forward price in the previous question is 128, describe
    the arbitrage strategy and compute the locked-in profit.
    \item A European call has $S_0=75$, $K=70$, $r=3\%$, $T=1$, and price 3.50.
    Does it violate the lower bound? Show the calculation.
    \item Given $S_0=100$, $K=105$, $r=5\%$, $T=1$, and $c=8$, use put-call
    parity to compute the European put price.
    \item A random variable is 10 with probability 0.25 and 20 with probability
    0.75. Compute its expectation and variance.
    \item Asset A has volatility 25\%, asset B has volatility 15\%, and their
    correlation is $-0.20$. Compute their covariance.
    \item A portfolio has weights 30\% and 70\% in the two assets from the
    previous question. Compute portfolio variance and volatility.
    \item Daily volatility is 0.9\%. Annualize it using 252 trading days.
    \item A regression beta is defined as
    $\beta=\Cov(R_i,R_m)/\Var(R_m)$. If $\Cov(R_i,R_m)=0.0012$ and
    $\Var(R_m)=0.0008$, compute beta.
    \item In a one-step binomial tree, $S_0=50$, $u=1.3$, $d=0.8$, $K=55$,
    $r=4\%$, and $T=1$. Price the European call by risk-neutral valuation.
    \item For the previous question, compute the replicating Delta.
    \item An option has Delta 0.40 and Gamma 0.03. Approximate the price change
    when the underlying rises by 5, ignoring time and volatility changes.
    \item A stock follows GBM with annual volatility 30\%. What is the
    approximate one-standard-deviation daily percentage shock using 252 trading
    days?
    \item Use It\^o's lemma to explain why $d\ln S$ has drift
    $\mu-\frac{1}{2}\sigma^2$ when $dS=\mu Sdt+\sigma Sdz$.
    \item A private equity fund has paid-in capital of 500, distributions of
    220, and residual value of 430. Compute DPI, RVPI, and TVPI.
    \item A project costs 1,000 and pays 400, 450, and 500 over the next three
    years. Compute NPV at a 9\% discount rate.
    \item Give one example of a real option and map it to a financial option
    payoff.
\end{enumerate}

\end{document}
